\documentclass{article}
\usepackage[T1]{fontenc}
\usepackage{hyperref}
\usepackage{xurl}
\usepackage{fancyhdr}
\usepackage[margin=0.8in]{geometry}

\newcommand{\projecttitle}{Unsupervised Grammar Induction for Magic: The Gathering Cards}
\newcommand{\studentname}{Avi Hyman}
\newcommand{\supervisorname}{Andrew Caines}
\newcommand{\dosname}{Timothy Jones}

\pagestyle{fancy}
\fancyhf{}
\fancyhead[C]{%
  \small
  \begin{tabular}{c}
    \textbf{\projecttitle}\\[2pt]
    \scriptsize Student: \studentname \quad | \quad Supervisor: \supervisorname \quad | \quad DoS: \dosname
  \end{tabular}%
}
\renewcommand{\headrulewidth}{0.4pt}
\setlength{\headheight}{38pt}
\setlength{\headsep}{8pt}

\begin{document}

\section{Introduction and Description of Work}
Magic: The Gathering (MTG) is a Collectible Card Game (CCG). Currently, over 30,000 unique MTG cards exist.

MTG cards' text is known as Oracle Text (OT). OT is designed to clarify ambiguity, interface with the comprehensive rulebook, and standardise the rules language.

The result of this is that OT lies somewhere between a natural language and an artificial language.

Academic papers have been written on MTG. Most notably, a 2019 paper proves that MTG is Turing Complete.

Previously, formal grammars have been successful in parsing a subset of OT. These tend to be hand-written rather than induced, likely due to a lack of pre-annotated data available in this field---meaning supervised grammar induction approaches are infeasible.\bigskip

I propose the production of an Unsupervised Grammar Induction (UGI) model for MTG cards, which learns a grammar for OT. This model will automatically parse OT into its inherent structure, which can lead to further applications including automated code generation and semantic search (outside the proposed scope of this project).

As an initial extension, I will analyse how OT compares to natural English, through evaluation of an MTG trained grammar on natural English, and an English trained grammar on OT. \bigskip

This project will involve these subtasks:
\begin{itemize}
    \item Research existing UGI approaches
    \begin{itemize}
        \item I will research each of the many approaches to UGI, determine which are most appropriate to my use-case, and decide upon one to implement algorithmically.
        \item I will research established evaluation methods---using both annotated and unannotated data---which fit my chosen approach.
    \end{itemize}
    \item Prepare test data
    \begin{itemize}
        \item I will extract card data from Scryfall API and perform filtering and seperation tasks, forming a dataset.
        \item I will annotate a subset of data by hand for the purpose of evaluation.
    \end{itemize}
    \item Implement data pre-processing
    \begin{itemize}
        \item I will implement steps to improve the UGI training quality, such as entity masking and lexing.
        \item I will split the per-card data into training-appropriate sentences.
    \end{itemize}
    \item Implement UGI training algorithm
    \begin{itemize}
        \item I will implement an algorithm in line with my research.
        \item I will run the model to produce an induced grammar for my dataset.
    \end{itemize}
    \item Evaluate performance
    \begin{itemize}
        \item I will evaluate against annotated and unannotated data, following the researched methods.
        \item I will evaluate performance of OT-trained grammar on a natural English dataset, and vice versa (initial extension)
    \end{itemize}
\end{itemize}

There is plenty of room for further extensions, the description of which is outside the scope of this proposal.

\section{Starting Point}

\begin{itemize}
        \item No work has been completed on any project code prior to the commencement date.
        \item I have no prior experience with Computational Linguistics outside of the Part IB Formal Models of Language course.
        \item I have completed some advance research on UGI methods at a high level of abstraction, primarily through this review paper: \url{https://doi.org/10.1017/S1351324920000327}
        \item I have previously used the Scryfall API to retrieve MTG data, but I have not used it in the sense of this project.
    \end{itemize}

\section{Success Criteria}

Project success depends on:
\begin{enumerate}
    \item A successful implementation of a UGI algorithm, capable of outputting a best-guess induced grammar given unannotated input data.
    \item A data gathering script which accesses an MTG API and stores a dataset locally. 
    \item A pre-processing script which accesses the dataset and outputs UGI-ready data.
    \item A script which evaluates the UGI output against hand-annotated and unannotated data, providing a metric comparable to those used in existing papers.
\end{enumerate}

\section{Plan and Milestones}

\subsection{12th--25th October}
\begin{itemize}
    \item Perform research to narrow down the potential UGI approaches to four candidates, with justification.
    \item Write program to gather the dataset (API access, initial filtering and formatting)
\end{itemize}

\subsection{26th October--8th November}
\begin{itemize}
    \item Perform deep research on two of the four candidates. Rate each approach in key heuristic areas.
    \item Start writing data pre-processing code---focus on approach-neutral processing. 
\end{itemize}

\subsection{9th--22nd November}
\begin{itemize}
    \item Perform deep research on the remaining two candidates. Rate each approach as before.
    \item Choose one approach to focus on.
    \item Continue writing pre-processing code.
\end{itemize}

\subsection{23rd November--6th December}
\begin{itemize}
    \item Research and plan algorithmic implementation of the chosen UGI approach.
    \item End of Michaelmas term---likely to be very busy period outside of the project.
\end{itemize}

\subsection{7th--20th December}
\begin{itemize}
    \item Begin implementation of the chosen UGI approach.
    \item Begin annotating data by hand for supervised evaluation.
\end{itemize}

\subsection{21st December--3rd January}
\begin{itemize}
    \item Complete implementation of the chosen UGI approach.
    \item Continue annotating data.
\end{itemize}

\subsection{4th--17th January}
\begin{itemize}
    \item Implement evaluation approaches.
    \item Run model end-to-end, obtain initial results.
\end{itemize}

\subsection{18th--31st January}
\begin{itemize}
    \item Buffer time for extensions.
    \item Write progress report.
\end{itemize}

\subsection{1st--14th February}
\begin{itemize}
    \item Continued buffer time for extensions.
    \item Project Report deadline: 5th February
\end{itemize}

\subsection{15th--28th February}
\begin{itemize}
    \item Final buffer time for extensions.
    \item Progress Report Presentation held by 26th February
\end{itemize}

\subsection{1st--14th March}
\begin{itemize}
    \item Begin dissertation writing.
    \item Write introduction chapter.
    \item Write preparation chapter.
\end{itemize}

\subsection{15th--28th March}
\begin{itemize}
    \item Begin implementation chapter.
\end{itemize}

\subsection{29th March--11th April}
\begin{itemize}
    \item Complete implementation chapter.
    \item Write evaluation chapter.
\end{itemize}

\subsection{12th--25th April}
\begin{itemize}
    \item Improve dissertation based on feedback.
\end{itemize}

\subsection{26th April--9th May}
\begin{itemize}
    \item Buffer time for dissertation writing.
\end{itemize}

\subsection{10th--14th May}
\begin{itemize}
    \item Finalise dissertation.
    \item Dissertation Deadline: 14th May
\end{itemize}

\section{Resource Declaration}
I will use my own Laptop - an MSI Sword 16 with an Intel i7-14700HX, 16GB RAM, and an RTX 4060 Laptop GPU. I accept full responsibility for this machine and I have made contingency plans to protect myself against hardware and/or software failure.

I will access the Scryfall API to gather my dataset. I will locally save the data gathered so as to avoid the risk inherent with relying upon a live API.

I will need to access a popular source of pre-annotated natural English data for evaluation for the initial extension. The exact source depends on the chosen UGI method. As above, I will store required data locally where possible.

I will use Git as a version control system with GitHub for backup.
\end{document}
